\subsection{Homomorphism and Isomorphism}

\begin{definition}
    Let $\left( G_1,* \right)$ and $\left( G_2,\circ \right)$
    be groups. A map $\varphi: G_1 \rightarrow G_2$ is called 
    a homomorphism if for all $a,b,c \in G_1$ we have 
    $\varphi\left( a*b \right)=\varphi\left( a \right)\circ \varphi\left( b \right)$\\
    If, an addition, $\varphi$ is a bijection we say $\varphi$ is an isomorphism
    and $G_1 \t{ and } G_2$ are isomorphic \[G_1\cong G_2\]
\end{definition}

\begin{theorem}
    \label[theorem]{thm:cyclicisomorphism}

    Suppose $G=\ang{g}$ is a cyclic group.\\
    If $|G|=\infty$ then $G \cong \Z$\\
    If $|G|=n$ then $G \cong \Z$
\end{theorem}
\begin{proof}
    Assume $|G|=\infty$ Define:
    \begin{align*}
        \varphi: G &\rightarrow \Z\\
        g^m &\mapsto m
    \end{align*}
    \begin{align*}
        \varphi^{-1}: \Z &\rightarrow G\\
        m &\mapsto g
    \end{align*}

    Let $g^a,g^b \in G$ and $c,d \in \Z$

    Consider: $\varphi(g^ag^b)=\varphi(g^{a+b})=a+b=\varphi(g^a)+\varphi(g^b)$
    
    Consider: $\varphi^{-1}(c+d)=g^{c+d}=g^cg^d=\varphi^{-1}(c)+\varphi^{-1}(d)$

    Consider: $\varphi\circ\varphi^{-1}(c)=\varphi(g^c)=c$

    $\varphi^{-1}\circ\varphi(g^a)=\varphi^{-1}(a)=g^a$

    Thus, $\varphi$ is a homomorphism and bijective

    $\implies G \cong \Z$\\\\

    Assume $|G|=n$
    \begin{align*}
        \varphi: G &\rightarrow \Z_n\\
        g^m &\mapsto [m]
    \end{align*}
    \begin{align*}
        \varphi^{-1}: \Z_n &\rightarrow G\\
        [m] &\mapsto g
    \end{align*}

    homomorphism: \checkmark

    Consider: $\varphi\circ\varphi^{-1}([c])=\varphi(g^c)=[c]$

    $\varphi^{-1}\circ\varphi(g^a)=\varphi^{-1}([a])=g^a$

    Thus, $\varphi$ is a homomorphism and bijective

    $\implies G \cong \Z_n$
\end{proof}

\begin{theorem}[Cayley's Theorem]
    \label[theorem]{thm:cayleysthm}

    Every group is isomorphic to a permutation group.
\end{theorem}

\begin{proof}
    Let $G$ be a group and by \Cref{lem:permutationgroups} $\implies S_G$ is a group\\

    Define:
    \begin{align*}
        \varphi: G &\rightarrow S_G\\
        g &\mapsto l_g\\
        l_g : G &\rightarrow G\\
        x &\mapsto gx
    \end{align*}

    \begin{itemize}
        \item Injective: Let $x,y \in G$ s.t. $l_g(x)=l_g(y)$
        
        $\implies gx=gy \implies x=y$

        \item Surjective: Let $y \in G$
        
        Consider: $l_g(g^{-1}y)=gg^{-1}y=y$
    \end{itemize}

    Thus $l_g \in S_G$

    Consider: $a,b \in G$

    $\varphi(ab)(x)=l_{ab}(x)=abx=a(bx)=l_a\circ l_b(x)=\varphi(a)\varphi(b)(x)$

    Injective: $a,b,x \in G$ s.t. $\varphi(a)=\varphi(b)$

    $l_a=l_b \implies l_a(x)=l_b(x)\;\; \forall x \in G \implies ax=bx \implies ae=be \implies a=b$

    $\implies G \cong \t{im(G)} \leq S_G$
\end{proof}
\newpage
\begin{theorem}
    \label[theorem]{thm:homomorphismgeneral}
    Suppose $\varphi: G_1 \rightarrow G_2$ is a homomorphism then
    \begin{enumerate}[I.]
        \item $\varphi(e_1)=e_2$
        \item $\forall n \in \Z, \varphi(g^n)=[\varphi(g)]^n$
        \item $ab=ba\implies \varphi(a)\varphi(b)=\varphi(b)\varphi(a)$
        \item ord$(\varphi(g))\mid\t{ord}(g)$
        \item $H \leq G_1 \implies \varphi(H)\leq G_2$
        \item $K \leq G_2 \implies \varphi^{-1}(K)\leq G_1$
    \end{enumerate}
\end{theorem}
\begin{enumerate}[I.]
    \item \begin{proof}
        $e_2\varphi(e_1)=e_2\varphi(e_1e_1)=e_2\varphi(e_1)\varphi(e_1)\implies
        e_2=\varphi(e_1)$
    \end{proof}
    \item \begin{proof}
        Base Case: $n=0\;\; \varphi(g^0)=\varphi(e_1)=e_2=[\varphi(g)]^0$

        Assume statement is true for $k\geq0$ for $k \in \N$

        Consider: $k+1$

        $\varphi(g^{k+1})=\varphi(g)\varphi(g^k)=\varphi(g)[\varphi(g)]^k=[\varphi(g)]^{k+1}$

        Note: $e_2=\varphi(e_1)=\varphi(gg^{-1})=\varphi(g)\varphi(g^{-1})\implies [\varphi(g)]^{-1}=\varphi(g^{-1})$

        This proves the statement for all negative integers since we just apply the induction on g then find its inverse
    \end{proof}
    \item \begin{proof}
        $\varphi(a)\varphi(b)=\varphi(ab)=\varphi(ba)=\varphi(b)\varphi(a)$
    \end{proof}
    \item \begin{proof}
        Set $n = \t{ord}(g)$

        $\implies \varphi(g)^n=\varphi(g^n)=\varphi(e_1)=e_2 \implies \t{ord}(\varphi(g)) \mid \t{ord}(g)$
    \end{proof}
    \item \begin{proof}
        Let $x,y \in \varphi(H)$ s.t. $\varphi(a)=x \t{ and } \varphi(b)=y$ for $a,b \in H$

        $\implies ab^{-1}\in H \implies \varphi(ab^{-1})=\varphi(a)\varphi(b)^{-1}\in \varphi(H)$

        Thus, by \Cref{thm:subgrouptest} $\varphi(H)\leq G_2$
    \end{proof}
    \item \begin{proof}
        Let $a,b \in \varphi^{-1}(K) \implies \varphi(a),\varphi(b)\in K$

        $\implies \varphi(ab^-1)=\varphi(a)\varphi(b)^{-1} \in K \implies ab^{-1}\in \varphi^{-1}(K)$

        Thus, by \Cref{thm:subgrouptest} $\varphi^{-1}(K)\leq G_1$
    \end{proof}
\end{enumerate}

\newpage
\begin{theorem}
    \label[theorem]]{thm:isomorphismgeneral}
    Suppose $\varphi: G_1 \rightarrow G_2$ is an isomorphism then
    \begin{enumerate}[I.]
        \item $\varphi^{-1}$ is an isomorphism
        \item $G_1$ is abelian $\iff$ $G_2$ is abelian
        \item $G_1$ is cyclic $\iff$ $G_2$ is cyclic
        \item ord$(g)=\t{ord}(\varphi(g))$
    \end{enumerate}
\end{theorem}
\begin{enumerate}[I.]
    \item \checkmark
    \item \begin{proof}
        Suppose $G_1$ is abelian and $x,y \in G_2$

        $xy=\varphi(a)\varphi(b)=\varphi(ab)=\varphi(ba)=\varphi(b)\varphi(a)=yx$

        Do same thing for other direction
    \end{proof}
    \item \begin{proof}
        Suppose $G_1$ is cyclic and $\ang{g}=G_1$

        Take $x \in G_2 \implies \varphi(a)=x$

        $a=g^n \implies x=\varphi(g^n)=\varphi(g)^n\implies x\in \ang{\varphi(g)}\implies G_2 \subseteq \ang{\varphi(g)}$

        Take $x \in \ang{\varphi(g)} \implies x=\varphi(g)^n$

        $\implies x=\varphi(g^n) \implies x\in G_2 \implies \ang{\varphi(g)} \subseteq G_2$

        $\implies G_2=\ang{\varphi(g)} \implies G_2$ is cyclic

        Do same thing for other direction
    \end{proof}
    \item \begin{proof}
        By \Cref{thm:homomorphismgeneral}.IV ord$(g) \mid \t{ord}(\varphi(g))$ and ord$(\varphi(g)) \mid \t{ord}(g) \implies \t{ord}(g) = \t{ord}(\varphi(g))$
    \end{proof}
\end{enumerate}

\newpage
\subsubsection{Exercises}
\begin{example}
    Describe all homomorphisms $\varphi : \Z_{12}\rightarrow\Z_{20}$

    ord$(1)=12 \t{ in } \Z_{12} \implies \t{ord}(\varphi(1)) \mid 12 \implies \t{ord}(\varphi(1))\in \left\{ 1,2,3,4,6,12 \right\}$

    ord$(m)=\frac{n}{\t(gcd)(m,n)}$ in $\Z_n$
    \begin{itemize}
        \item ord$(\varphi(1))=1 \implies \varphi(1)=0\implies \varphi(n)=0$
        \item ord$(\varphi(1))=2 \implies \varphi(1)=10$ 
        
        this is the only element with a gcd of 10 with 20 $\implies \varphi(n)=10n$

        \item ord$(\varphi(1))=3,6,12$ there is no element with these orders since these are not divsors of 10
        \item ord$(\varphi(1))=4 \implies \varphi(1)=5,15 \implies \varphi(n)=5n,15n$
    \end{itemize}
\end{example}

\begin{example}
    Show that $\R \ncong \R^\times$

    Assume $\R \cong \R^\times$ via $\varphi: \R^\times \rightarrow \R$

    Take $x \in R^\times$ s.t. $\varphi(x)=-1$

    Consider $\varphi(\frac{x}{2})^2=\varphi(x)=-1 \implies \varphi(\frac{x}{2})\in \R \t{ with } \varphi(\frac{x}{2})^2=-1 \contr$
\end{example}